\documentclass[11pt,letterpaper]{article}
\usepackage[margin=1in]{geometry}
\usepackage{lmodern}
\usepackage[T1]{fontenc}
\usepackage{amsmath,amssymb,amsthm}
\usepackage{booktabs}
\usepackage{microtype}
\usepackage{hyperref}

\newtheorem{theorem}{Theorem}
\newtheorem{lemma}[theorem]{Lemma}
\newtheorem{proposition}[theorem]{Proposition}
\theoremstyle{definition}
\newtheorem{convention}[theorem]{Convention}

\title{PAPER\_S205 --- Session 205 Backfill: \\
       Phase H205 $E^{\pm}(t)$ Expansion--Erosion Footprint \\
       (Universal Buoyancy Counter-Balance, $F_{U,Bi}/F_U$)}
\author{Daniel T. Murphy}
\date{May 18, 2026}

\begin{document}
\maketitle

\begin{abstract}
Session 205 introduced three standalone modules: \texttt{positive\_et\_expansion.py}
(positive-energy expansion engine $E^{+}(t)$, Kozima neutron-drop coupling,
expansion Lagrangian), \texttt{negative\_et\_erosion.py} (negative-energy
erosion engine $E^{-}(t)$, net-energy balance, gravitational-wave damping,
erosion Lagrangian), and \texttt{uqff\_vs\_string\_comparison.py} (10-aspect
weighted UQFF vs. String-Theory comparison). The master equations encode the
universal-buoyancy counter-balance $F_{U,Bi}/F_U$ as the single algebraic
parameter that governs whether vacuum energy expands, erodes, or balances.
This paper records the six exact identities (the H205 closures) that pin
that algebra in the closure ledger. Three of the six are closed-form
calculus / algebra theorems (H205-1 doubling-time, H205-2 critical balance,
H205-3 net-factor); two are definitional arithmetic checkpoints (H205-4
weight-sum, H205-6 test count); one is the postulated VDS dimensionality
(H205-5, $N_{\text{levels}}=26$). No claim is made on the overall winner of
the UQFF vs. String comparison; the closures only verify that the
comparison's scoring weights are a probability simplex.
\end{abstract}

\section{Ledger Summary}

\begin{center}
\begin{tabular}{lllc}
\toprule
ID & Closure & Value & Class \\
\midrule
H205-1 & Doubling-time identity $t_{\text{double}}(\kappa + [\mathrm{SSq}]/26) = \ln 2$ & $31.6172$~s & THEOREM \\
H205-2 & Critical balance ratio $r_{\text{crit}}$ solving $2r-1=0$ & $1/2$ & THEOREM \\
H205-3 & Net-factor formula $\text{nf}(r) = 2r - 1$, $r=1.1$ & $1.2$ & THEOREM \\
H205-4 & Comparison-weight unit sum & $1.00$ & arithmetic \\
H205-5 & VDS spacetime level count $N_{\text{levels}}$ & $26$ & CONVENTION \\
H205-6 & S205 module-test total $9+10+14$ & $33$ & arithmetic \\
\bottomrule
\end{tabular}
\end{center}

All six register as \texttt{EXACT} in \texttt{master\_closures.csv} /
\texttt{sigma\_table.csv} (script: \texttt{\_session205\_closures.py}; final
parseable line \texttt{S205-doubling-time-identity: 31.6172 vs 31.6172 -> EXACT}).
Independent verification is provided by \texttt{\_verify\_s205.py}.

\section{Theorems --- Calculus and Algebraic Identities}

\begin{theorem}[H205-1 --- Doubling-time identity]
\label{thm:H205-1}
Let $E^{+}(t) = E_0 \exp\!\big((\kappa + [\mathrm{SSq}]/N_{\mathrm{lev}})\, t\big)
\cdot S_{26}([\mathrm{SSq}]) \cdot (F_{U,Bi}/F_U)$, with all factors other than
$\exp(\cdot)$ time-independent. Then the doubling time
$t_{\mathrm{double}}$ satisfying $E^{+}(t_{\mathrm{double}}) = 2\, E^{+}(0)$ is
\[
   t_{\mathrm{double}} \;=\; \frac{\ln 2}{\kappa + [\mathrm{SSq}]/N_{\mathrm{lev}}}.
\]
\end{theorem}
\begin{proof}
$E^{+}(t)/E^{+}(0) = \exp\!\big((\kappa + [\mathrm{SSq}]/N_{\mathrm{lev}})\, t\big)$
since the $S_{26}$ and buoyancy-ratio factors cancel. Setting this equal to $2$
and taking $\ln$ gives the stated identity.
\end{proof}

\noindent
With canonical $\kappa = 5 \times 10^{-4}\,\text{day}^{-1} = 5.787 \times 10^{-9}\,\text{s}^{-1}$,
$[\mathrm{SSq}]=0.57$, $N_{\mathrm{lev}}=26$:
\[
   t_{\mathrm{double}} = \frac{\ln 2}{5.787\times10^{-9} + 0.57/26}
   = \frac{0.6931\ldots}{2.1923\ldots \times 10^{-2}} = 31.6172\ldots\,\text{s},
\]
matching the module test $T_5$ output \texttt{3.1617e+01 s} byte-for-byte.

\begin{theorem}[H205-2 --- Critical balance ratio]
\label{thm:H205-2}
The net-energy equation
\[
   E_{\mathrm{net}}(t) \;=\; E_0\, \exp\!\big((\kappa + [\mathrm{SSq}]/N_{\mathrm{lev}})\, t\big)\, S_{26}([\mathrm{SSq}])\, \big[\, 2(F_{U,Bi}/F_U) - 1\, \big]
\]
admits the unique zero $E_{\mathrm{net}}(t)\equiv 0$ at $r := F_{U,Bi}/F_U = 1/2$.
\end{theorem}
\begin{proof}
The exponential and $S_{26}$ factors are strictly positive. Hence
$E_{\mathrm{net}}=0$ iff $2r - 1 = 0$, which gives the unique root $r=1/2$.
\end{proof}

\begin{theorem}[H205-3 --- Net-factor formula]
\label{thm:H205-3}
Define $\text{nf}(r) := r - (1-r) = 2r - 1$, where $r = F_{U,Bi}/F_U$.
Then for $r = 1.1$, $\text{nf}(1.1) = 1.2$ exactly.
\end{theorem}
\begin{proof}
$\text{nf}(1.1) = 2(1.1) - 1 = 2.2 - 1 = 1.2$.
\end{proof}

\section{Structural Checkpoints --- Arithmetic and Convention}

\begin{proposition}[H205-4 --- Weight-sum identity]
\label{prop:H205-4}
The category weights
$W = \{w_{\text{test}}, w_{\text{pred}}, w_{\text{found}}, w_{\text{math}}\}
= \{0.30, 0.30, 0.20, 0.20\}$ used in
\texttt{uqff\_vs\_string\_comparison.OverallScoreCalculator} satisfy
$\sum_{w \in W} w = 1$, i.e.\ $W$ is a probability simplex.
\end{proposition}
\begin{proof}
$0.30 + 0.30 + 0.20 + 0.20 = 1.00$ by direct addition.
\end{proof}

\begin{convention}[H205-5 --- VDS spacetime level count]
\label{conv:H205-5}
The Vacuum-Dimensional-Spacetime (VDS) level count is fixed at
$N_{\mathrm{lev}} = 26$ throughout the UQFF framework. This appears in the
$S_{26}([\mathrm{SSq}])$ factor of $E^{\pm}(t)$ and in the per-level rate
contribution $[\mathrm{SSq}]/N_{\mathrm{lev}}$. It is a postulated structural
constant, not a derived quantity in S205.
\end{convention}

\begin{proposition}[H205-6 --- S205 module-test total]
\label{prop:H205-6}
The three S205 standalone modules carry, respectively, $9$, $10$, and $14$
self-test cases (\texttt{positive\_et\_expansion.py},
\texttt{negative\_et\_erosion.py}, \texttt{uqff\_vs\_string\_comparison.py}).
Their sum is $33$, and all $33$ pass clean.
\end{proposition}
\begin{proof}
$9 + 10 + 14 = 33$ by direct addition; module self-tests verified
independently by execution.
\end{proof}

\section{Honest Classification}

Of the six H205 closures: three are closed-form theorems
(\ref{thm:H205-1}, \ref{thm:H205-2}, \ref{thm:H205-3}); two are definitional
arithmetic checkpoints (\ref{prop:H205-4}, \ref{prop:H205-6}); one is a
postulated structural constant (\ref{conv:H205-5}). No physical-content claim
is made; in particular, the $E^{\pm}(t)$ framework's identification of
$F_{U,Bi}/F_U = 1/2$ as the cosmological balance point is a model statement
of the UQFF Universal-Buoyancy postulate, not a derivation from first
principles. The closures verify only that the model's algebra is internally
consistent.

\section{Reproducibility}

\begin{itemize}
\item Closure script: \texttt{\_session205\_closures.py}
\item Closure JSON:   \texttt{\_session205\_closures.json}
\item Tier-13 block:  \texttt{UQFF\_UNIFIED\_CLOSURE\_DERIVATIONS.py},
                      records \texttt{H205-1} through \texttt{H205-6}
\item Ledger row:     \texttt{master\_closures.csv} (S205 row 6, status \texttt{OK})
\item Sigma row:      \texttt{sigma\_table.csv}     (S205 row 6, bucket \texttt{EXACT})
\item Verifier:       \texttt{\_verify\_s205.py} (independent algorithm:
                      \texttt{sympy} symbolic root, \texttt{fractions.Fraction}
                      rational arithmetic, \texttt{mpmath} arbitrary-precision
                      logarithm)
\item Source modules: \texttt{positive\_et\_expansion.py},
                      \texttt{negative\_et\_erosion.py},
                      \texttt{uqff\_vs\_string\_comparison.py}
\item Canonical record: \texttt{VALIDATION\_MASTER\_INDEX\_2.md}, L462 (v5.62,
                        commit \texttt{336d98e7}).
\end{itemize}

\end{document}
